% TOP_PROBLEM: {"problemId": "problem.bqp-versus-np", "canonicalTitle": "BQP Versus NP: Complete Language-Class Relation", "releaseRank": 16, "primaryDomain": "theoretical_computer_science", "primaryDomainLabel": "Theoretical computer science", "displayStatus": "open", "releaseStatus": "open"}
% !TeX program = lualatex
% Definition and sources only; this file is not an LLM solution attempt.
\documentclass[11pt]{article}
\usepackage[margin=25mm]{geometry}
\usepackage{fontspec}
\setmainfont{Latin Modern Roman}
\usepackage{amsmath,amssymb,mathtools}
\usepackage{unicode-math}
\setmathfont{Latin Modern Math}
\usepackage{xurl,hyperref}
\hypersetup{colorlinks=true,urlcolor=blue}
\setcounter{secnumdepth}{0}
\setlength{\parindent}{0pt}
\setlength{\parskip}{5pt}
\setlength{\emergencystretch}{3em}
\begin{document}
\section*{\texorpdfstring{BQP Versus NP: Complete Language-Class Relation}{BQP Versus NP: Complete Language-Class Relation}}\label{16}
\subsection{Definitions and mathematical statement}
A total language belongs to $\mathsf{BQP}$ if uniformly generated polynomial-size quantum circuits accept yes-inputs with probability at least $2/3$ and no-inputs with probability at most $1/3$. Use precisely the fixed gates, ancillas, and measurement convention in the recorded scope. The class $\mathsf{NP}$ instead uses polynomial-length classical witnesses and deterministic verification. Determine the pair of truth values
\[
 \bigl(\mathsf{NP}\subseteq\mathsf{BQP},\quad
       \mathsf{BQP}\subseteq\mathsf{NP}\bigr).
\]
A complete answer distinguishes equality, one of the two strict containments, or incomparability. Oracle access, advice, and arbitrary uncomputable gate constants are excluded.
\par\smallskip\textit{Formulation and concept references:} \href{https://arxiv.org/pdf/0804.3401v1}{[S1]}.

\subsection{Short English statement}
Which is stronger: efficient quantum decision-making or efficient classical verification, or are their powers incomparable?
\subsection{Sources}
\begin{itemize}
\item[S1]John Watrous, Quantum Computational Complexity, arXiv:0804.3401v1; language questions corroborated by Scott Aaronson lecture notes.
\url{https://arxiv.org/pdf/0804.3401v1}.
\textit{Formulation source.}
\item[S2]Complexity Theory, Lecture 25: Quantum Computing (2) — Markus Krötzsch, January 27, 2025.
\url{https://iccl.inf.tu-dresden.de/w/images/4/4d/CT2024-Lecture-25-overlay.pdf}.
\textit{Further reference; not a proof endorsement.}
\item[S3]Introduction to Quantum Information Science lecture notes — Scott Aaronson.
\url{https://www.scottaaronson.com/qclec.pdf}.
\textit{Further reference; not a proof endorsement.}
\item[S4]Random Permutations in Computational Complexity — Hitchcock, Sekoni and Shafei (November 11, 2025).
\url{https://arxiv.org/abs/2511.08786}.
\textit{Further reference; not a proof endorsement.}
\end{itemize}
\end{document}
